%=======================================================================
%  LABORATORIO N.o 6 - EL TRANSISTOR EN CORTE Y SATURACION
%  Curso : Circuitos Electronicos  (2026-2)
%  UCSP  - Facultad de Ingenieria - Ing. Electronica y de Telecomunicaciones
%  Docente     : Ebert San Roman Castillo
%  Laboratorio : Ing. Jeanpier Ancori Sanchez
%
%  Version 2 - Pedido del docente: que el alumno CALCULE los valores del
%  interruptor a partir de condiciones dadas. Junta en una sola practica la
%  version 1 (rehecha desde "Lab. 05 - Aplicaciones con el transistor en
%  corte y saturacion") con la "Parte 1. Diseno de circuito en corte y
%  saturacion" del "Lab. 06 - Diseno con transistor bipolar BJT" del docente
%  (2026-1: V_BB = V_CC = 10 V, I_C(sat) = 30 mA, LED azul).
%  Tres experimentos, todos con la misma estructura (Datos -> Calcule ->
%  Arme y mida -> Anote) y un cuadro resumen al inicio del desarrollo:
%    1. Diseno del interruptor de un LED azul y prueba "al limite" con
%       R_B = 47 k en BC548B y 2N2222A (reemplaza al potenciometro de la v1).
%    2. Compuertas NOR y NAND con diodos (las de la v1). El inversor NOT de
%       la v1 sale: el experimento 1 ya es un inversor.
%    3. Rele de 5 V con diodo volante.
%  Cambios frente al original del docente: I_C(sat) = 20 mA y no 30 mA (30 mA
%  es el maximo absoluto de un LED de 5 mm y con 12 V obligaria a R_C de
%  1 W); V_CC = 12 V y V_BB = 5 V por los rieles de la Analog Discovery
%  Studio. La parte 2 del original (divisor de tension) la cubre el Lab. 5.
%  Formato: cuestionario de 3 preguntas, conclusiones 5 + 5; no se imprimen
%  los resultados.
%
%  VALORES DE REFERENCIA PARA EL DOCENTE (a mano: V_BE = 0,7 V,
%  V_CE(sat) = 0,2 V, LED azul 3,0 V, LED rojo 2,0 V a 12 mA y 1,85 V a
%  2 mA, diodo 0,65 V; no se imprimen):
%    Exp. 1: R_C = (12 - 3 - 0,2)/20 mA = 440 -> 470 Ohm (E12, hacia arriba),
%      I_C(sat) = 18,7 mA, P_RC = 165 mW (0,5 W). I_B = 1,87 mA,
%      R_B = 4,3/1,87 mA = 2,30 k -> 2,2 k (hacia abajo), I_B = 1,95 mA,
%      beta forzado = 9,6. V_BB = 0: corte; con el multimetro (10 MOhm) el
%      LED azul cae ~2 V y se mide V_C ~ 10 V.
%      Al limite con h_FE(min) = 200: R_B = 200*4,3/18,7 mA = 46 k -> 47 k.
%    Exp. 2 NOR (diodos + R_B = 1 k + 22 k a tierra):
%      00 : corte, D = 2,6-2,7 V, LED encendido.
%      01, 10, 11 : nodo X = 4,35 V, I_B ~ 3,6 mA, saturacion, D ~ 0,1 V.
%    Exp. 2 NAND (22 k de pull-up, D3 en serie con la base):
%      00, 01, 10 : X ~ 0,6 V, no alcanza para D3 + V_BE (~1,3 V): corte,
%                   E = 2,6-2,7 V, LED encendido.
%      11 : X ~ 1,35 V, I_B = (5 - 1,35)/22 k = 0,17 mA, I_C(sat) = 4,8 mA,
%           beta forzado = 29, saturacion, E ~ 0,1-0,2 V.
%      Cuestionario 2: R_P(max) = beta (5 - 1,4)/4,8 mA = 112 k (beta 150).
%      Entrada al aire: en la NAND cuenta como 1 (el diodo no conduce); en la
%      NOR cuenta como 0 (la 22 k lleva el nodo a tierra).
%    Exp. 3 (rele SRD-05VDC-SL-C, ~70 Ohm, 2N2222A, I_B <= 4 mA):
%      I_C(sat) = 4,8/70 = 68,6 mA ; con beta forzado 20, I_B = 3,43 mA,
%      R_B = 1,25 k -> 1,2 k, I_B = 3,58 mA, beta forzado 19,2
%      (h_FE(min) = 50 a 150 mA). Al abrir, con 1N4004: V_C = 5,7 V.
%      Cuestionario 3: 0,1 H * 68,6 mA / 1 us = 6,9 kV >> V_CEO = 40 V.
%    Cuestionario 1: con h_FE(min) = 75 del 2N2222A, R_B = 17 k -> 15 k,
%      beta forzado = 65.
%
%  SIMULACION (ngspice 26 de Fusion/Eagle; modelos BC548B de NXP, 2N2222A,
%  1N4004 y LED azul y rojo ajustados a mano; netlists, datos y graficas en
%  ../Simulaciones):
%    Exp. 1, BC548B, R_B = 2,2 k: V_BE 0,746 V, V_CE 0,059 V, I_B 1,93 mA,
%            I_C 19,0 mA, V_LED 2,99 V. 2N2222A: V_CE 0,049 V.
%            R_B = 47 k : BC548B V_CE = 0,17 V (satura justo, beta ~ 207)
%                         2N2222A V_CE = 1,42 V (activa, beta ~ 178) <-- clave
%    Exp. 2 NOR : 00 -> D 2,65 V ; 01, 10, 11 -> X 4,36 V, I_B 3,59 mA,
%                 D 0,015 V.
%           NAND: 00, 01, 10 -> X 0,46-0,49 V, V_BE 0,44-0,46 V, E 2,65 V ;
%                 11 -> X 1,18 V, I_B 0,174 mA, I_C 4,94 mA, E 0,056 V.
%    Exp. 3, L = 0,1 H supuesta: I_C 69,8 mA, I_B 3,50 mA, V_CE 0,118 V.
%            Pico al abrir: 5,79 V con 1N4004; 3,4 kV sin diodo (el modelo
%            no tiene ruptura; en la realidad el transistor se dana).
%=======================================================================

\documentclass[11pt,a4paper]{article}

%------------------------------------------------------------- idioma
% Compilar con XeLaTeX (no con pdflatex): Trebuchet MS es una fuente TrueType
% del sistema y solo fontspec, bajo XeLaTeX o LuaLaTeX, sabe cargarla.
\usepackage{fontspec}
\setmainfont{Trebuchet MS}[
  Path            = C:/Windows/Fonts/,
  Extension       = .ttf,
  UprightFont     = trebuc,
  BoldFont        = trebucbd,
  ItalicFont      = trebucit,
  BoldItalicFont  = trebucbi,
  Ligatures       = TeX
]
\usepackage[spanish,es-nodecimaldot]{babel}
\usepackage{microtype}

%----------------------------------------------------------- geometria
\usepackage[left=2.4cm,right=2.4cm,top=3.1cm,bottom=2.6cm,headheight=30pt,
            footskip=1.2cm]{geometry}

%------------------------------------------------------------ graficos
\usepackage{graphicx}
\usepackage{xcolor}
\usepackage{float}
\usepackage{placeins}
\usepackage[font=small,labelfont=bf,labelsep=period]{caption}
\usepackage{subcaption}
\graphicspath{{imagenes/}}
\usepackage[american,siunitx]{circuitikz}
% babel-spanish activa < > " y rompe las opciones de TikZ sin esta libreria.
\usetikzlibrary{calc,patterns,babel}

%-------------------------------------------------------------- tablas
\usepackage{array}
\usepackage{tabularx}
\usepackage{multirow}

%----------------------------------------------------- matematicas / SI
\usepackage{amsmath,amssymb}
\usepackage{siunitx}
\sisetup{
  output-decimal-marker = {,},
  per-mode              = symbol,
  mode                  = match,
  detect-weight         = true,
  detect-family         = true
}

%--------------------------------------------------------------- listas
\usepackage{enumitem}
\setlist[itemize]{itemsep=3pt,topsep=4pt,parsep=0pt}
\setlist[enumerate]{itemsep=8pt,topsep=4pt,parsep=0pt}

%------------------------------------------------------------ comandos
\usepackage{pgffor}

% Lineas punteadas para respuestas del alumno: \lineasresp{n}
\newcommand{\lineasresp}[1]{%
  \par\vspace{6pt}%
  \foreach \i in {1,...,#1}{\noindent\rule{\linewidth}{0.4pt}\par\vspace{8pt}}%
}

% Recuadro "Figura / Comentario": una celda alta para pegar la captura del
% osciloscopio y otra debajo para el comentario.
\newcommand{\cuadrorespuesta}{%
  \par\vspace{6pt}\noindent
  \begin{tabularx}{\linewidth}{|>{\centering\arraybackslash}X|}
    \hline
    \rule{0pt}{4.4cm}\raisebox{2.1cm}{\textcolor{black!45}{Figura}}\\
    \hline
    \rule{0pt}{2.0cm}\raisebox{0.9cm}{\textcolor{black!45}{Comentario}}\\
    \hline
  \end{tabularx}\par\vspace{8pt}%
}

% Celda vacia de altura comoda para escribir a mano.
\newcommand{\hueco}{\rule{0pt}{0.5cm}}

% Fila de una tabla de verdad: entradas fijas y el resto en blanco.
% \filaverdad{A}{B} para dos entradas.
\newcommand{\filaverdad}[2]{#1 & #2 & \hueco & & & & \\ \hline}

% Tabla de verdad de una compuerta de dos entradas.
% #1 = nombre de la salida, #2 = leyenda, #3 = etiqueta.
\newcommand{\tablaverdad}[3]{%
  \begin{table}[H]
    \centering
    \renewcommand{\arraystretch}{1.15}
    \begin{tabularx}{\linewidth}{|c|c|*{3}{>{\centering\arraybackslash}X|}>{\centering\arraybackslash}m{2.0cm}|>{\centering\arraybackslash}m{2.4cm}|}
      \hline
      \multicolumn{2}{|c|}{\textbf{Entradas}}
        & \multicolumn{3}{c|}{\textbf{Salida $V_{#1}$}}
        & \multirow{2}{=}{\centering\textbf{LED}\par (on/off)}
        & \multirow{2}{=}{\centering\textbf{Estado de Q}\par (corte/sat.)} \\
      \cline{1-5}
      $A$ & $B$ & \textbf{Teórico} & \textbf{Simulado} & \textbf{Medido} & & \\
      \hline
      \filaverdad{\qty{0}{\volt}}{\qty{0}{\volt}}
      \filaverdad{\qty{0}{\volt}}{\qty{5}{\volt}}
      \filaverdad{\qty{5}{\volt}}{\qty{0}{\volt}}
      \filaverdad{\qty{5}{\volt}}{\qty{5}{\volt}}
    \end{tabularx}
    \caption{#2}
    \label{#3}
  \end{table}%
}


%----------------------------------------------------- compaginacion
\usepackage{needspace}
% \needspace mide mal el espacio libre despues de figuras y tablas [H] y deja
% paginas medio vacias; \Needspace es la variante exacta del mismo paquete.
\let\needspace\Needspace
\usepackage{etoolbox}
\pretocmd{\section}{\needspace{5\baselineskip}}{}{}
\pretocmd{\subsection}{\needspace{4\baselineskip}}{}{}
\widowpenalty=10000
\clubpenalty=10000

%---------------------------------------------------- encabezado / pie
\usepackage{fancyhdr}
\pagestyle{fancy}
\fancyhf{}
\fancyhead[L]{\footnotesize UCSP --- Facultad de Ingeniería\\
              Ing. Electrónica y de Telecomunicaciones}
\fancyhead[R]{\footnotesize 2026-2 Circuitos Electrónicos\\
              Ebert San Román Castillo}
\fancyfoot[L]{\footnotesize Laboratorio N.\textsuperscript{o} 6 --- El transistor en corte y saturación}
\fancyfoot[R]{\footnotesize Página \thepage\ de \pageref{LastPage}}
\renewcommand{\headrulewidth}{0.4pt}
\renewcommand{\footrulewidth}{0.4pt}
\usepackage{lastpage}

%------------------------------------------------------- numeracion
\renewcommand{\thesection}{\Roman{section}}
\renewcommand{\thesubsection}{\arabic{subsection}}
\renewcommand{\thesubsubsection}{\thesubsection.\arabic{subsubsection}}
\setcounter{secnumdepth}{3}

%-------------------------------------------------------- hiperenlaces
\usepackage[hidelinks,pdfusetitle]{hyperref}
\usepackage{xurl}
\hypersetup{
  pdftitle={Laboratorio 6 - El transistor en corte y saturacion},
  pdfsubject={Circuitos Electronicos - Guia de practicas},
  pdfauthor={Ebert San Roman Castillo; Jeanpier Ancori Sanchez},
  pdfkeywords={electronica, transistor, BJT, corte, saturacion, interruptor,
               diseno, beta forzado, compuertas logicas, NOT, NOR, NAND, rele, diodo volante, TTL, LED,
               Analog Discovery Studio, Proteus, Multisim, osciloscopio}
}

% babel-spanish rotula las tablas como "Cuadro"; los originales dicen "Tabla".
\addto\captionsspanish{\renewcommand{\tablename}{Tabla}}

\begin{document}

%=======================================================================
%  CABECERA DE LA GUIA
%=======================================================================
\begin{center}
    {\normalsize\bfseries SEGUNDA UNIDAD: TRANSISTOR DE UNIÓN BIPOLAR (BJT)}
\end{center}

\vspace{0.55cm}

\begingroup
\renewcommand{\tabularxcolumn}[1]{m{#1}}
\renewcommand{\arraystretch}{2.0}
\setlength{\tabcolsep}{12pt}
\noindent
\begin{tabularx}{\textwidth}{|>{\centering\arraybackslash}m{2.2cm}|>{\centering\arraybackslash}X|>{\centering\arraybackslash}m{3.0cm}|}
\hline
\rule{0pt}{2.1cm}{\fontsize{58}{58}\selectfont\bfseries 6}
    & \begin{minipage}[c]{\linewidth}\centering
          {\large\bfseries Guía de Prácticas}\\[7pt]
          {\Large\bfseries El Transistor en\\ Corte y Saturación}
      \end{minipage}
    & {\Large\bfseries Nota}\\
\hline
\multicolumn{3}{|l|}{\rule{0pt}{1.3cm}\large\textbf{Grupo:}~\rule{4.6cm}{0.4pt}%
    \hfill\large\textbf{Fecha:}~\rule{4.6cm}{0.4pt}\hspace{2mm}}\\
\hline
\multicolumn{3}{|l|}{\rule{0pt}{1.2cm}\large\textbf{Alumno(s):}~\rule{11.2cm}{0.4pt}\hspace{2mm}}\\
\multicolumn{3}{|l|}{\rule{0pt}{1.2cm}\phantom{\large\textbf{Alumno(s):}}~\rule{11.2cm}{0.4pt}\hspace{2mm}}\\
\hline
\end{tabularx}
\endgroup

\vspace{0.8cm}

% Encabezado de cada paso dentro de un experimento: \paso{1}{Datos}
\newcommand{\paso}[2]{%
  \par\needspace{10\baselineskip}\vspace{8pt}\noindent
  \colorbox{black!8}{\makebox[\dimexpr\linewidth-2\fboxsep\relax][l]{%
    \textbf{Paso #1 \,·\, #2}}}\par\vspace{4pt}}

%=======================================================================
\section{Objetivos}
%=======================================================================
\begin{itemize}
    \item Calcular las resistencias de base y de colector de un transistor bipolar que trabaja
          como interruptor, a partir de las condiciones de la carga y de la hoja de datos, y
          comprobar por simulación y por medición que queda en corte o en saturación.
    \item Explicar por qué una resistencia de base calculada con el $\beta$ justo puede no
          saturar a otro transistor.
    \item Verificar compuertas NOR y NAND con diodos y un transistor, y manejar un relé
          protegiendo al transistor con un diodo volante.
\end{itemize}

%=======================================================================
\section{Contenido teórico}
%=======================================================================
\begin{itemize}
    \item Regiones de trabajo del transistor de unión bipolar (BJT): corte, activa y
          saturación.
    \item El transistor como interruptor: corriente de saturación y $\beta$ forzado.
    \item Compuertas lógicas con diodos y transistores (lógica DTL); cargas inductivas y
          diodo volante.
\end{itemize}

%=======================================================================
\section{Equipos y materiales}
%=======================================================================
\noindent\textbf{Laboratorio:} Electrónica de Comunicaciones.

\vspace{0.25cm}
\noindent\textbf{Equipos y dispositivos:}
\begin{itemize}
    \item Tarjeta Analog Discovery Studio: rieles fijos de \qty{+12}{\volt} y
          \qty{+5}{\volt}, generador de funciones \texttt{W1} y osciloscopio de dos canales.
\end{itemize}

\noindent\textbf{Software:}
\begin{itemize}
    \item Simulador Multisim o Proteus.
\end{itemize}

\vspace{0.25cm}
\needspace{6\baselineskip}
\noindent\textbf{Materiales y fungibles} (deben ser adquiridos por los alumnos):
\begin{itemize}
    % --- Semiconductores
    \item 02 transistores de unión bipolar NPN BC548B y 02 2N2222A (uno de repuesto de cada
          uno). Sus terminales están en distinto orden: revise la hoja de datos antes de armar.
    \item 02 LED azules y 04 LED rojos de \qty{5}{\milli\metre} (uno de repuesto de cada color).
    \item 05 diodos rectificadores 1N4004 (uno de repuesto).
    \item 01 relé de \qty{5}{\volt} con un contacto normalmente abierto (tipo
          SRD-05VDC-SL-C, bobina de unos \qty{70}{\ohm}).
    % --- Resistencias (la que mas disipa es RC = 470 Ohm del exp. 1: 165 mW)
    \item Resistencias a \qty{0.5}{\watt}: 02 de \qty{470}{\ohm}.
    \item Resistencias a \qty{0.25}{\watt}:
          02 de \qty{2.2}{\kilo\ohm},
          01 de \qty{47}{\kilo\ohm},
          04 de \qty{1}{\kilo\ohm},
          02 de \qty{1.2}{\kilo\ohm},
          03 de \qty{330}{\ohm} y
          03 de \qty{22}{\kilo\ohm}.
          Traiga además algunos valores vecinos (\qty{1.8}{\kilo\ohm},
          \qty{2.7}{\kilo\ohm}, \qty{1.5}{\kilo\ohm}) por si su diseño previo da un valor
          distinto.
    % --- Montaje
    \item Protoboard.
    \item Jumpers macho-macho para protoboard.
    \item Alicates para cortar y pelar alambre.
    \item Multímetro (con prueba de diodos).
    \item Lentes de protección y mandil.
\end{itemize}

\FloatBarrier
%=======================================================================
\section{Conocimientos previos}
%=======================================================================

%-----------------------------------------------------------------------
\subsection{El transistor como interruptor}
%-----------------------------------------------------------------------
Un transistor que conmuta trabaja solo en dos estados: en \textbf{corte} no circula
corriente de colector y es un \textbf{interruptor abierto}; en \textbf{saturación}
$V_{CE} \approx \qty{0.2}{\volt}$ y es un \textbf{interruptor cerrado}
(figura~\ref{fig:generico}).

\begin{figure}[H]
  \centering
  \begin{circuitikz}[scale=0.85, transform shape]
    \draw (4.2,0) node[npn](Q){};
    \coordinate (c)   at ($(Q.C)+(0,0.5)$);
    \coordinate (top) at ($(c)+(0,3.4)$);
    \draw (top) to[R, l=$R_C$] ++(0,-1.7) to[generic, l=carga] (c) -- (Q.C);
    \draw (top) node[circ]{} -- ++(1.6,0) node[right]{$V_{CC}$};
    \draw ($(Q.B)+(-3.0,0)$) node[left]{$V_{BB}$}
          to[R, l=$R_B$, o-, i>^=$I_B$] (Q.B);
    \draw (Q.C) ++(0,0.25) node[right, font=\small]{$\downarrow I_C$};
    \draw (Q.E) node[ground]{};
  \end{circuitikz}
  \caption{Interruptor con transistor NPN. La carga puede ser un LED o la bobina de un relé.}
  \label{fig:generico}
\end{figure}

\noindent\textbf{Cómo se diseña.} Se parte de lo que la carga necesita:
\begin{enumerate}[label=\arabic*), leftmargin=*, itemsep=4pt]
  \item \textbf{Corriente de saturación} y resistencia de colector:
        \begin{equation}
          R_C = \frac{V_{CC} - V_{\text{carga}} - V_{CE(\text{sat})}}{I_{C(\text{sat})}}
          \label{eq:rc}
        \end{equation}
        Si la carga es la bobina de un relé no hay $R_C$:
        $I_{C(\text{sat})} = (V_{CC} - V_{CE(\text{sat})})/R_{\text{bobina}}$.
  \item \textbf{Corriente de base}, con un $\beta$ forzado mucho menor que el $\beta$ del
        transistor (la regla práctica es 10):
        \begin{equation}
          I_B = \frac{I_{C(\text{sat})}}{\beta_{\text{forzado}}}
          \label{eq:ib}
        \end{equation}
  \item \textbf{Resistencia de base:}
        \begin{equation}
          R_B = \frac{V_{BB} - V_{BE}}{I_B}
          \label{eq:rb}
        \end{equation}
  \item \textbf{Valores comerciales (E12):} $R_C$ se redondea \textbf{hacia arriba} y $R_B$
        \textbf{hacia abajo}. Con ellos se recalculan $I_{C(\text{sat})}$, $I_B$ y
        $\beta_{\text{forzado}} = I_{C(\text{sat})}/I_B$, y se revisa la potencia
        $P = I^2R$ de cada resistencia.
\end{enumerate}

\noindent\textbf{¿Por qué $\beta_{\text{forzado}} = 10$?} El $\beta$ cambia mucho de una
unidad a otra: la hoja de datos solo garantiza un mínimo ($h_{FE(\min)}$). Si $R_B$ se
calcula con el $\beta$ justo,
\begin{equation}
  R_{B(\max)} = \frac{\beta\,(V_{BB} - V_{BE})}{I_{C(\text{sat})}},
  \label{eq:rbmax}
\end{equation}
un transistor con un $\beta$ menor no se satura: queda en la región activa, $V_{CE}$ sube y
la carga recibe menos corriente.

\needspace{14\baselineskip}
\noindent\textbf{Ejemplo resuelto} (valores distintos a los de la práctica). LED rojo de
\qty{2}{\volt} con $I_{C(\text{sat})} = \qty{15}{\milli\ampere}$, $V_{CC} = \qty{9}{\volt}$ y
$V_{BB} = \qty{3.3}{\volt}$.
\begin{table}[H]
  \centering
  \renewcommand{\arraystretch}{2.1}
  \begin{tabularx}{\linewidth}{|>{\centering\arraybackslash}m{2.6cm}|>{\centering\arraybackslash}X|>{\centering\arraybackslash}m{2.8cm}|}
    \hline
    \textbf{Paso} & \textbf{Cálculo} & \textbf{E12} \\ \hline
    $R_C$ & $\dfrac{9 - 2 - \num{0.2}}{\qty{15}{\milli\ampere}} = \qty{453}{\ohm}$
          & \qty{470}{\ohm} (arriba) \\ \hline
    $I_{C(\text{sat})}$, $I_B$ & $\dfrac{\num{6.8}}{\qty{470}{\ohm}} = \qty{14.5}{\milli\ampere}$,
            \qquad $\dfrac{\qty{14.5}{\milli\ampere}}{10} = \qty{1.45}{\milli\ampere}$ & --- \\ \hline
    $R_B$ & $\dfrac{\num{3.3} - \num{0.7}}{\qty{1.45}{\milli\ampere}} = \qty{1.79}{\kilo\ohm}$
          & \qty{1.5}{\kilo\ohm} (abajo) \\ \hline
    Comprobación & $I_B = \dfrac{\num{2.6}}{\qty{1.5}{\kilo\ohm}} = \qty{1.73}{\milli\ampere}$,
                   \qquad $\beta_{\text{forzado}} = \dfrac{\num{14.5}}{\num{1.73}} = \num{8.4}$
                 & $\leq 10$: satura \\ \hline
  \end{tabularx}
  \caption{Ejemplo de diseño de un interruptor para un LED.}
  \label{tab:ejemplo}
\end{table}

%-----------------------------------------------------------------------
\subsection{¿Corte, activa o saturación?}
\label{sec:identificar}
%-----------------------------------------------------------------------
\noindent\textbf{Por cálculo.} Calcule $I_B$ siguiendo el camino de la base (reste
\qty{0.7}{\volt} por cada diodo que conduce). Si no hay corriente de base, el transistor está
en \textbf{corte}. Si la hay, compare $\beta\,I_B$ con $I_{C(\text{sat})}$: si
$\beta\,I_B \geq I_{C(\text{sat})}$, está en \textbf{saturación}; si no, está en la
\textbf{región activa}.

\noindent\textbf{Por medición.} Mida $V_{BE}$ y $V_{CE}$ y use la tabla~\ref{tab:criterios}.

\begin{table}[H]
  \centering
  \renewcommand{\arraystretch}{1.3}
  \begin{tabularx}{\linewidth}{|>{\centering\arraybackslash}m{2.1cm}|*{3}{>{\centering\arraybackslash}X|}}
    \hline
    \textbf{Región} & $\boldsymbol{V_{BE}}$ & $\boldsymbol{V_{CE}}$ & \textbf{El interruptor está} \\
    \hline
    Corte      & $< \qty{0.5}{\volt}$ & $\approx V_{CC}$ & abierto \\ \hline
    Activa     & \qtyrange{0.6}{0.7}{\volt} & $> \qty{0.3}{\volt}$ & a medias (no sirve) \\ \hline
    Saturación & \qtyrange{0.7}{0.8}{\volt} & $\leq \qty{0.3}{\volt}$ & cerrado \\ \hline
  \end{tabularx}
  \caption{Cómo identificar la región de trabajo a partir de lo que se mide.}
  \label{tab:criterios}
\end{table}

%-----------------------------------------------------------------------
\subsection{Compuertas lógicas con diodos y un transistor}
%-----------------------------------------------------------------------
Con \qty{0}{\volt} como \textbf{0 lógico} y \qty{5}{\volt} como \textbf{1 lógico}, un
transistor que conmuta es un \textbf{inversor}: con la entrada en 1 se satura y su colector
queda cerca de \qty{0}{\volt}; con la entrada en 0 se corta y su colector sube. Delante del
inversor se pone una red de diodos:
\begin{itemize}
  \item \textbf{Red OR} (figura~\ref{fig:nor}): basta con que una entrada esté en 1 para que
        su diodo conduzca y llegue corriente a la base. Con el inversor forma una \textbf{NOR}.
  \item \textbf{Red AND} (figura~\ref{fig:nand}): basta con que una entrada esté en 0 para que
        su diodo conduzca y le quite la corriente a la base. Con el inversor forma una
        \textbf{NAND}.
\end{itemize}
Un LED en la salida indica el nivel: encendido es 1 y apagado es 0.

%-----------------------------------------------------------------------
\subsection{Cargas inductivas y diodo volante}
%-----------------------------------------------------------------------
La bobina de un relé se opone a que su corriente cambie. Cuando el transistor se corta, la
bobina genera un voltaje $v_L = L\,di/dt$ que puede llegar a cientos de voltios y dañar al
transistor. El \textbf{diodo volante}, en paralelo con la bobina y con el cátodo hacia
$V_{CC}$, le da a esa corriente un camino y limita el colector a
$V_{CC} + \qty{0.7}{\volt}$.

\FloatBarrier
\needspace{14\baselineskip}
%=======================================================================
\section{Desarrollo del laboratorio}
%=======================================================================
Cada experimento tiene los mismos cuatro pasos. Los pasos~1 y~2 (\textbf{Datos} y
\textbf{Calcule}) se traen resueltos de casa, junto con la simulación en Multisim o Proteus;
los pasos~3 y~4 (\textbf{Arme y mida} y \textbf{Anote}) se hacen en el laboratorio. Use
$V_{BE} = \qty{0.7}{\volt}$ y $V_{CE(\text{sat})} = \qty{0.2}{\volt}$.

\begin{table}[H]
  \centering
  \renewcommand{\arraystretch}{1.35}
  \begin{tabularx}{\linewidth}{|>{\centering\arraybackslash}m{1.0cm}|>{\raggedright\arraybackslash}m{3.0cm}|>{\raggedright\arraybackslash}X|>{\raggedright\arraybackslash}X|>{\centering\arraybackslash}m{1.3cm}|}
    \hline
    \textbf{Exp.} & \textbf{Circuito} & \textbf{En casa: calcule} & \textbf{En el laboratorio: mida}
      & \textbf{Riel} \\ \hline
    1 & Interruptor de un LED azul
      & $R_C$, $R_B$ y $\beta_{\text{forzado}}$; $R_{B(\max)}$ con $h_{FE(\min)}$
      & $V_{BE}$, $V_{CE}$, $I_B$, $I_C$ con la entrada en 1 y en 0; $V_{CE}$ con
        $R_B = \qty{47}{\kilo\ohm}$ en los dos transistores
      & \qty{+12}{\volt} \\ \hline
    2 & Compuertas NOR y NAND
      & Estado del transistor y salida para cada combinación de entradas
      & La salida, el nodo $X$ y $V_{BE}$ en cada combinación
      & \qty{+5}{\volt} \\ \hline
    3 & Relé con diodo volante
      & $I_{C(\text{sat})}$, $R_B$ y $\beta_{\text{forzado}}$
      & $V_{CE}$, $I_B$, $I_C$; el pico del colector al abrir el relé
      & \qty{+5}{\volt} \\ \hline
  \end{tabularx}
  \caption{Resumen de la práctica.}
  \label{tab:resumen}
\end{table}

%-----------------------------------------------------------------------
\subsection{Experimento 1: interruptor de un LED azul}
%-----------------------------------------------------------------------
Una salida lógica de \qty{5}{\volt} debe encender un LED azul alimentado con \qty{12}{\volt}.

\paso{1}{Datos}
\begin{table}[H]
  \centering
  \renewcommand{\arraystretch}{1.2}
  \begin{tabularx}{0.8\linewidth}{|>{\raggedright\arraybackslash}X|>{\centering\arraybackslash}m{4.4cm}|}
    \hline
    $V_{CC}$ & \qty{12}{\volt} \\ \hline
    $V_{BB}$ (entrada $A$) & \qty{5}{\volt} (1) y \qty{0}{\volt} (0) \\ \hline
    Corriente del LED, $I_{C(\text{sat})}$ & \qty{20}{\milli\ampere} \\ \hline
    $V_{LED}$ (LED azul) & \qty{3.0}{\volt} \\ \hline
    $\beta_{\text{forzado}}$ & 10 \\ \hline
    Transistor & BC548B ($h_{FE(\min)} = 200$) \\ \hline
  \end{tabularx}
  \caption{Experimento 1: datos de diseño.}
  \label{tab:cond1}
\end{table}

\begin{figure}[H]
  \centering
  \begin{circuitikz}[scale=0.9, transform shape]
    \draw (4.2,0) node[npn](Q){};
    \node[right] at (Q.east) {BC548B};
    \coordinate (c)   at ($(Q.C)+(0,0.5)$);
    \coordinate (l)   at ($(c)+(0,1.9)$);
    \coordinate (top) at ($(l)+(0,2.0)$);
    \draw (top) to[R, l=$R_C$, a=?] (l)
                to[leDo, l=LED azul] (c) -- (Q.C);
    \draw (top) node[circ]{} -- ++(1.8,0) node[right]{$V_{CC} = \qty{+12}{\volt}$};
    \draw (c) to[short, *-o] ++(2.0,0) node[right]{$C$};
    \draw ($(Q.B)+(-3.0,0)$) node[left]{$A$}
          to[R, l=$R_B$, a=?, o-] (Q.B);
    \draw (Q.E) node[ground]{};
  \end{circuitikz}
  \caption{Experimento 1: interruptor de un LED azul.}
  \label{fig:led}
\end{figure}

\paso{2}{Calcule (en casa)}
\begin{enumerate}[label=\alph*), leftmargin=*, itemsep=3pt]
  \item $R_C$ con~\eqref{eq:rc}, $I_B$ con~\eqref{eq:ib} y $R_B$ con~\eqref{eq:rb}, como en el
        ejemplo de la tabla~\ref{tab:ejemplo}. Elija los valores E12 y recalcule
        $I_{C(\text{sat})}$, $I_B$, $\beta_{\text{forzado}}$ y la potencia en $R_C$.
  \item $R_{B(\max)}$ con~\eqref{eq:rbmax} y $\beta = 200$. Debe salir cerca de
        \qty{47}{\kilo\ohm}.
  \item Simule el circuito con sus valores, con la entrada en \qty{5}{\volt} y en
        \qty{0}{\volt}. Simule también $R_B = \qty{47}{\kilo\ohm}$ con el BC548B y con el
        2N2222A.
\end{enumerate}

\paso{3}{Arme y mida (en el laboratorio)}
\begin{enumerate}[label=\alph*), ref=\alph*, leftmargin=*, itemsep=3pt]
  \item Arme el circuito con sus valores E12 y conecte la tarjeta como en la
        figura~\ref{fig:conexion}.
  \item Ponga \texttt{W1} en \qty{5}{\volt} continuos y luego en \qty{0}{\volt}. En cada caso
        mida $V_{BE}$, $V_{CE}$ y las caídas en $R_B$ y $R_C$; calcule $I_B$ e $I_C$ con la ley
        de Ohm.
  \item \label{paso:cuadrada}Ponga \texttt{W1} como onda cuadrada de \qty{1}{\kilo\hertz},
        amplitud \qty{2.5}{\volt} y offset \qty{2.5}{\volt}, y capture $v_A$ (CH1) y $v_C$
        (CH2) con marcadores en los dos niveles de $v_C$.
  \item Con \texttt{W1} en \qty{5}{\volt}, cambie $R_B$ por \qty{47}{\kilo\ohm} y mida
        $V_{CE}$. Repita con el 2N2222A en lugar del BC548B (cuide el orden de sus
        terminales) y compare el brillo del LED.
\end{enumerate}

% Conexion con la tarjeta, en los colores del arnes de la Analog Discovery
% Studio (W1 amarillo, osciloscopio 1 naranja y 2 azul, GND negro,
% alimentacion roja). Los canales son diferenciales: 1- y 2- van a GND.
\begin{figure}[H]
  \centering
  \begin{circuitikz}[scale=0.85, transform shape]
    \colorlet{cw1}{yellow!70!black}
    \colorlet{cch1}{orange!90!black}
    \colorlet{cch2}{blue!75!black}
    \colorlet{cvcc}{red!80!black}
    \coordinate (gnd) at (0,-1.3);
    \draw (4.2,0) node[npn](Q){};
    \node[right] at (Q.east) {BC548B};
    \coordinate (c)   at ($(Q.C)+(0,0.5)$);
    \coordinate (l)   at ($(c)+(0,1.7)$);
    \coordinate (top) at ($(l)+(0,1.8)$);
    \draw (top) to[R, l=$R_C$] (l)
                to[leDo] (c) -- (Q.C);
    \coordinate (a) at (0.9,0);
    \draw (a) to[R, l=$R_B$, *-] (Q.B);
    \draw (Q.E) -- (Q.E |- gnd) node[circ]{};
    % bus de tierra
    \coordinate (p2m) at (6.3,0);
    \coordinate (p2p) at (7.1,0);
    \coordinate (p5)  at (7.9,0);
    \draw[very thick] (-0.8,-1.3) -- (p2m |- gnd);
    % tarjeta
    \coordinate (ytab) at (0,-2.9);
    \draw[thick, fill=black!6, rounded corners=3pt]
         (-1.6,-2.9) rectangle (8.8,-4.4);
    \node[font=\bfseries] at (3.6,-3.95) {Analog Discovery Studio};
    \foreach \x/\n/\c in {-0.8/GND/black, 0.3/W1/cw1, 0.9/1+/cch1, 1.5/1$-$/cch1}{
      \fill[\c] (\x-0.12,-3.02) rectangle (\x+0.12,-2.78);
      \node[font=\scriptsize\ttfamily, below] at (\x,-3.02) {\n};
    }
    \foreach \p/\n/\c in {p2m/2$-$/cch2, p2p/2+/cch2, p5/+12V/cvcc}{
      \fill[\c] ($(\p |- ytab)+(-0.12,-0.12)$) rectangle ($(\p |- ytab)+(0.12,0.12)$);
      \node[font=\scriptsize\ttfamily, below] at ($(\p |- ytab)+(0,-0.12)$) {\n};
    }
    % cables
    \draw[black, line width=1.3pt] (-0.8,-2.78) -- (-0.8,-1.3) node[circ]{};
    \draw[cw1, line width=1.3pt]  (0.3,-2.78) -- (0.3,0) -- (a);
    \draw[cch1, line width=1.3pt] (0.9,-2.78) -- (a);
    \draw[cch1, line width=1.3pt, dashed] (1.5,-2.78) -- (1.5,-1.3) node[circ, black]{};
    \draw[cch2, line width=1.3pt, dashed] ($(p2m |- ytab)+(0,0.12)$) -- (p2m |- gnd) node[circ, black]{};
    \draw[cch2, line width=1.3pt] ($(p2p |- ytab)+(0,0.12)$) -- (p2p |- c) -- (c) node[circ, black]{};
    \draw[cvcc, line width=1.3pt] ($(p5 |- ytab)+(0,0.12)$) -- (p5 |- top) -- (top);
    % rotulos
    \node[cch1, font=\scriptsize, above] at ($(a)+(0,0.1)$) {CH1: $v_A$};
    \node[cch2, font=\scriptsize, above] at ($(c)!0.5!(p2p |- c)$) {CH2: $v_C$};
  \end{circuitikz}
  \caption{Conexión con la Analog Discovery Studio. \texttt{W1} da la entrada $A$; \texttt{1+}
           la mide y \texttt{2+} mide el colector. \texttt{1$-$} y \texttt{2$-$} van a
           \texttt{GND}. El riel de \qty{+12}{\volt} alimenta el circuito. En el experimento 3
           se usa la misma conexión con el riel de \qty{+5}{\volt}.}
  \label{fig:conexion}
\end{figure}

\paso{4}{Anote}
\begin{table}[H]
  \centering
  \renewcommand{\arraystretch}{1.15}
  \begin{tabularx}{\linewidth}{|>{\centering\arraybackslash}m{3.0cm}|*{6}{>{\centering\arraybackslash}X|}}
    \hline
    & \multicolumn{3}{c|}{\textbf{Entrada en \qty{5}{\volt}}}
    & \multicolumn{3}{c|}{\textbf{Entrada en \qty{0}{\volt}}} \\
    \cline{2-7}
    \textbf{Parámetro} & Teórico & Simulado & Medido & Teórico & Simulado & Medido \\
    \hline
    $R_C$ y $R_B$ (E12)          & \hueco & \multicolumn{5}{c|}{} \\ \hline
    $V_{BE}$                     & \hueco & & & & & \\ \hline
    $V_{CE}$                     & \hueco & & & & & \\ \hline
    $I_B$                        & \hueco & & & & & \\ \hline
    $I_C$                        & \hueco & & & & & \\ \hline
    $\beta_{\text{forzado}} = I_C/I_B$ & \hueco & & & \multicolumn{3}{c|}{---} \\ \hline
    Estado de Q (corte/sat.)     & \hueco & & & & & \\ \hline
  \end{tabularx}
  \caption{Experimento 1: el interruptor diseñado.}
  \label{tab:exp1}
\end{table}

\begin{table}[H]
  \centering
  \renewcommand{\arraystretch}{1.15}
  \begin{tabularx}{0.85\linewidth}{|>{\centering\arraybackslash}m{3.2cm}|*{3}{>{\centering\arraybackslash}X|}}
    \hline
    $\boldsymbol{R_B = \qty{47}{\kilo\ohm}}$ & \textbf{$V_{CE}$ simulado} & \textbf{$V_{CE}$ medido}
      & \textbf{Estado de Q} \\ \hline
    BC548B  & \hueco & & \\ \hline
    2N2222A & \hueco & & \\ \hline
  \end{tabularx}
  \caption{Experimento 1: la misma $R_B$ en dos transistores distintos.}
  \label{tab:limite}
\end{table}

\noindent\begin{minipage}{\linewidth}
\noindent\textbf{Captura de la conmutación} (paso~3\ref{paso:cuadrada}).
\cuadrorespuesta
\end{minipage}\par

%-----------------------------------------------------------------------
\subsection{Experimento 2: compuertas NOR y NAND}
%-----------------------------------------------------------------------
Dos compuertas lógicas hechas con diodos y un transistor que solo trabaja en corte o en
saturación.

\paso{1}{Datos}
$V_{CC} = \qty{5}{\volt}$; entradas en \qty{0}{\volt} (0) o \qty{5}{\volt} (1); transistor
2N2222A con $\beta = 150$; diodos de \qty{0.7}{\volt}; LED rojo de \qty{2}{\volt}. Los
circuitos son los de las figuras~\ref{fig:nor} y~\ref{fig:nand}.

\begin{figure}[H]
  \centering
  \begin{circuitikz}[scale=0.85, transform shape]
    \draw (6.4,0) node[npn](Q){};
    \node[right] at ($(Q.E)+(0.05,0.25)$) {2N2222A};
    \coordinate (gnd) at ($(Q.E)+(0,-1.6)$);
    \coordinate (c)   at ($(Q.C)+(0,0.5)$);
    \coordinate (top) at ($(c)+(0,2.2)$);
    \draw (top) to[R, l=$R_C$, a=\qty{1}{\kilo\ohm}] (c) -- (Q.C);
    \draw (top) node[circ]{} -- ++(3.6,0) node[right]{$V_{CC} = \qty{+5}{\volt}$};
    \draw (c) to[short, *-] ++(2.9,0) coordinate(o)
          to[R, l=$R_L$, a=\qty{330}{\ohm}, *-] ($(o)+(0,-1.9)$)
          to[leDo] (o |- gnd);
    \draw (o) to[short, -o] ++(1.6,0) node[right]{$D$};
    \draw (Q.E) -- (Q.E |- gnd) node[circ]{} -- (o |- gnd);
    \coordinate (x) at ($(Q.B)+(-2.6,0)$);
    \draw ($(x)+(-3.0,1.2)$) node[left]{$A$} to[D, l=$D_1$, o-] ($(x)+(0,1.2)$) -- (x);
    \draw ($(x)+(-3.0,0)$) node[left]{$B$} to[D, l=$D_2$, o-*] (x);
    \draw (x) to[R, l=$R_B$, a=\qty{1}{\kilo\ohm}] (Q.B);
    \draw (x) to[R, l_=$R_P$, a^=\qty{22}{\kilo\ohm}] (x |- gnd) -- (Q.E |- gnd);
    \draw (x |- gnd) -- ++(0,-0.01) node[ground]{};
    \node[below right, font=\scriptsize] at (x) {$X$};
  \end{circuitikz}
  \caption{Compuerta NOR: red OR de diodos seguida de un inversor.}
  \label{fig:nor}
\end{figure}

\begin{figure}[H]
  \centering
  \begin{circuitikz}[scale=0.85, transform shape]
    \draw (6.4,0) node[npn](Q){};
    \node[right] at ($(Q.E)+(0.05,0.25)$) {2N2222A};
    \coordinate (gnd) at ($(Q.E)+(0,-1.6)$);
    \coordinate (c)   at ($(Q.C)+(0,0.5)$);
    \coordinate (top) at ($(c)+(0,2.2)$);
    \draw (top) to[R, l=$R_C$, a=\qty{1}{\kilo\ohm}] (c) -- (Q.C);
    \draw (c) to[short, *-] ++(2.9,0) coordinate(o)
          to[R, l=$R_L$, a=\qty{330}{\ohm}, *-] ($(o)+(0,-1.9)$)
          to[leDo] (o |- gnd);
    \draw (o) to[short, -o] ++(1.6,0) node[right]{$E$};
    \draw (Q.E) -- (Q.E |- gnd) -- (o |- gnd);
    \node[ground] at (Q.E |- gnd) {};
    \coordinate (x) at ($(Q.B)+(-2.4,0)$);
    \draw (x |- top) to[R, l=$R_P$, a=\qty{22}{\kilo\ohm}] (x);
    \draw (x |- top) -- (top) node[circ]{} -- ++(3.6,0) node[right]{$V_{CC} = \qty{+5}{\volt}$};
    \draw (x) to[D, l=$D_3$, *-] (Q.B);
    \draw (x) to[D, l_=$D_1$, -o] ($(x)+(-3.0,0)$) node[left]{$A$};
    \draw (x) -- ($(x)+(0,-1.2)$) to[D, l_=$D_2$, -o] ($(x)+(-3.0,-1.2)$) node[left]{$B$};
    \node[above left, font=\scriptsize] at (x) {$X$};
  \end{circuitikz}
  \caption{Compuerta NAND: red AND de diodos seguida de un inversor. $D_3$ sube en unos
           \qty{0.7}{\volt} el voltaje que necesita el nodo $X$ para saturar al transistor.}
  \label{fig:nand}
\end{figure}

\paso{2}{Calcule (en casa)}
\begin{enumerate}[label=\alph*), leftmargin=*, itemsep=3pt]
  \item Para cada compuerta y cada combinación de entradas, decida si el transistor está en
        corte o en saturación (sección~\ref{sec:identificar}) y calcule el voltaje de salida.
        Cuando el transistor se satura, sus \qty{0.2}{\volt} no encienden el LED de salida.
  \item Simule las dos compuertas con las cuatro combinaciones.
\end{enumerate}

\paso{3}{Arme y mida (en el laboratorio)}
\begin{enumerate}[label=\alph*), ref=\alph*, leftmargin=*, itemsep=3pt]
  \item Arme la NOR y aliméntela con el riel de \qty{+5}{\volt}. Para poner una entrada en 1,
        conéctela con un jumper al riel; para ponerla en 0, a \texttt{GND}. No deje entradas
        al aire.
  \item En cada combinación mida la salida, el nodo $X$ y $V_{BE}$, y observe el LED.
  \item Repita con la NAND.
\end{enumerate}

\paso{4}{Anote}
\tablaverdad{D}{Experimento 2: tabla de verdad de la compuerta NOR.}{tab:nor}
\tablaverdad{E}{Experimento 2: tabla de verdad de la compuerta NAND.}{tab:nand}

%-----------------------------------------------------------------------
\subsection{Experimento 3: manejo de un relé}
%-----------------------------------------------------------------------
Una salida de \qty{5}{\volt} que solo puede entregar \qty{4}{\milli\ampere}, como el pin de
un microcontrolador, debe activar un relé. El contacto del relé enciende un LED.

\paso{1}{Datos}
\begin{table}[H]
  \centering
  \renewcommand{\arraystretch}{1.2}
  \begin{tabularx}{0.8\linewidth}{|>{\raggedright\arraybackslash}X|>{\centering\arraybackslash}m{4.4cm}|}
    \hline
    $V_{CC}$ y $V_{BB}$ & \qty{5}{\volt} \\ \hline
    $R_{\text{bobina}}$ & mídala (unos \qty{70}{\ohm}) \\ \hline
    $I_B$ máxima que da la entrada & \qty{4}{\milli\ampere} \\ \hline
    $\beta_{\text{forzado}}$ & 20 \\ \hline
    Transistor & 2N2222A ($h_{FE(\min)} = 50$) \\ \hline
  \end{tabularx}
  \caption{Experimento 3: datos de diseño.}
  \label{tab:cond3}
\end{table}

\begin{figure}[H]
  \centering
  \begin{circuitikz}[scale=0.9, transform shape]
    \draw (4.2,0) node[npn](Q){};
    \node[right] at ($(Q.E)+(0.05,0.25)$) {2N2222A};
    \coordinate (c)   at ($(Q.C)+(0,0.5)$);
    \coordinate (top) at ($(c)+(0,3.0)$);
    \draw (top) to[L, l_=Bobina del relé] (c) -- (Q.C);
    \draw (c) to[short, *-] ++(2.4,0) coordinate(d)
          to[D, l_=$D_V$] (d |- top) -- (top);
    \node[right, font=\small] at ($(d)!0.5!(d |- top)+(0.55,-0.45)$) {1N4004};
    \draw (top) node[circ]{} -- ++(-1.8,0) node[left]{$V_{CC} = \qty{+5}{\volt}$};
    \draw (c) to[short, -o] ++(-1.1,0) node[left]{$C$};
    \draw ($(Q.B)+(-3.0,0)$) node[left]{$A$}
          to[R, l=$R_B$, a=?, o-] (Q.B);
    \draw (Q.E) node[ground]{};
    % contacto del rele
    \coordinate (k) at (10.2,3.5);
    \draw (k) node[circ]{} node[above]{$V_{CC}$}
          to[nos, l=contacto NA] ++(0,-1.6)
          to[R, l=\qty{330}{\ohm}] ++(0,-1.8)
          to[leDo, l=LED rojo] ++(0,-1.6) node[ground]{};
  \end{circuitikz}
  \caption{Experimento 3: relé con diodo volante $D_V$. El contacto normalmente abierto (NA),
           que la bobina mueve, enciende el LED rojo. \textbf{No retire $D_V$ con el circuito
           energizado.}}
  \label{fig:rele}
\end{figure}

\paso{2}{Calcule (en casa)}
\begin{enumerate}[label=\alph*), leftmargin=*, itemsep=3pt]
  \item $I_{C(\text{sat})}$ con la resistencia de la bobina, $I_B$ con~\eqref{eq:ib} y
        $\beta_{\text{forzado}} = 20$, y $R_B$ con~\eqref{eq:rb}. Elija el valor E12 y
        compruebe que $I_B$ no pase de \qty{4}{\milli\ampere}.
  \item Simule el circuito con y \textbf{sin} el diodo volante, y anote el voltaje pico del
        colector cuando el transistor se corta.
\end{enumerate}

\paso{3}{Arme y mida (en el laboratorio)}
\begin{enumerate}[label=\alph*), ref=\alph*, leftmargin=*, itemsep=3pt]
  \item Arme el circuito con el riel de \qty{+5}{\volt}. El cátodo de $D_V$ (la franja) va
        hacia $V_{CC}$.
  \item Ponga \texttt{W1} en \qty{5}{\volt} continuos y luego en \qty{0}{\volt}. Mida
        $V_{BE}$, $V_{CE}$ y las caídas en $R_B$ y en la bobina; calcule $I_B$ e $I_C$.
  \item \label{paso:volante}Ponga \texttt{W1} como onda cuadrada de \qty{2}{\hertz}, amplitud
        \qty{2.5}{\volt} y offset \qty{2.5}{\volt}. Con CH1 en $A$ y CH2 en el colector,
        dispare con el flanco de bajada de CH1 a \qty{1}{\milli\second} por división y capture
        el pico de $v_C$. \textbf{No retire el diodo para medir}: ese caso solo se simula.
\end{enumerate}

\paso{4}{Anote}
\begin{table}[H]
  \centering
  \renewcommand{\arraystretch}{1.15}
  \begin{tabularx}{\linewidth}{|>{\centering\arraybackslash}m{3.0cm}|*{6}{>{\centering\arraybackslash}X|}}
    \hline
    & \multicolumn{3}{c|}{\textbf{Entrada en \qty{5}{\volt}}}
    & \multicolumn{3}{c|}{\textbf{Entrada en \qty{0}{\volt}}} \\
    \cline{2-7}
    \textbf{Parámetro} & Teórico & Simulado & Medido & Teórico & Simulado & Medido \\
    \hline
    $R_{\text{bobina}}$ y $R_B$ (E12) & \hueco & \multicolumn{5}{c|}{} \\ \hline
    $V_{CE}$                     & \hueco & & & & & \\ \hline
    $I_B$                        & \hueco & & & & & \\ \hline
    $I_C$                        & \hueco & & & & & \\ \hline
    LED rojo (on/off)            & \hueco & & & & & \\ \hline
    Estado de Q (corte/sat.)     & \hueco & & & & & \\ \hline
    Pico de $v_C$ al abrir       & \multicolumn{3}{l|}{Con $D_V$:} & \multicolumn{3}{l|}{Sin $D_V$ (simulado):} \\ \hline
  \end{tabularx}
  \caption{Experimento 3: el relé.}
  \label{tab:exp3}
\end{table}

\noindent\begin{minipage}{\linewidth}
\noindent\textbf{Captura de la apertura del relé} (paso~3\ref{paso:volante}).
\cuadrorespuesta
\end{minipage}\par

\FloatBarrier
%=======================================================================
\section{Informe final}
%=======================================================================
\begin{enumerate}[label=\textbf{\arabic*.}, leftmargin=*, labelsep=0.6em]
  \item Adjunte los cálculos y las simulaciones de los tres experimentos.
  \item Compare los valores teóricos, simulados y medidos y explique las diferencias más
        grandes.
  \item Explique con la tabla~\ref{tab:limite} por qué la misma $R_B$ satura a un transistor
        y no al otro.
\end{enumerate}

\FloatBarrier
%=======================================================================
\section{Cuestionario}
%=======================================================================
\begin{enumerate}[label=\textbf{\arabic*.}, leftmargin=*, labelsep=0.6em]
\item En el experimento 1, ¿cuál es la mayor $R_B$ que satura a \emph{cualquier} 2N2222A, si su
$h_{FE(\min)}$ a esa corriente es 75? ¿Cuánto vale $\beta_{\text{forzado}}$ con esa $R_B$?
\lineasresp{3}

\item En la compuerta NAND, ¿cuál es el mayor valor de $R_P$ con el que el transistor todavía
se satura con las dos entradas en 1? ¿Por qué una entrada al aire cuenta como 1 en la NAND y
como 0 en la NOR?
\lineasresp{3}

\item En el experimento 3, ¿por qué el colector llega a unos \qty{5.7}{\volt} al abrir el relé
con el diodo volante? Sin el diodo, si la corriente de la bobina cae a cero en
\qty{1}{\micro\second} y la bobina tiene \qty{0.1}{\henry}, ¿qué voltaje aparecería?
Compárelo con el $V_{CEO}$ del 2N2222A (\qty{40}{\volt}).
\lineasresp{3}
\end{enumerate}

\FloatBarrier
%=======================================================================
\section{Conclusiones}
%=======================================================================
\begin{itemize}
    \item Redacte como mínimo 05 conclusiones \textbf{cualitativas} y 05
          \textbf{cuantitativas}.
\end{itemize}

\noindent\textbf{Cualitativas}
\lineasresp{10}

\needspace{6\baselineskip}
\noindent\textbf{Cuantitativas}
\lineasresp{10}

\end{document}
